% Part: model-theory
% Chapter: models-of-arithmetic
% Section: computable-models

\documentclass[../../../include/open-logic-section]{subfiles}

\begin{document}

\olfileid[pt]{mod}{mar}{cmp}
\section{Modelos computáveis da aritmética}

\begin{explain}
O modelo padrão~$\Struct{N}$ tem duas características convenientes.
Seu domínio é o conjunto dos números naturais~$\Nat$, O.s.\@, seus
elementos são justamente o tipo de coisa de que queremos falar usando
a linguagem da aritmética, e o numeral padrão~$\num{n}$ de fato
designa~$n$. A outra característica conveniente é que as interpretações
dos símbolos não lógicos de~$\Lang{L_A}$ são todas \emph{computáveis}.
As funções sucessor, adição e multiplicação que servem como
$\Assign{\prime}{N}$, $\Assign{+}{N}$ e $\Assign{\times}{N}$ são
funções computáveis de números. (Computáveis por máquinas de Turing ou
definíveis por recursão primitiva, por exemplo.) E a relação menor que
em~$\Struct{N}$, O.s.\@, $\Assign{<}{N}$, é decidível.

Modelos não padrão de teorias aritméticas como $\Th{Q}$ e $\Th{PA}$
devem conter elementos não padrão. Assim, seus domínios em geral
incluem !!{element}s além de~$\Nat$. No entanto, qualquer !!{structure}
enumerável pode ser construída sobre qualquer conjunto !!{denumerable},
inclusive~$\Nat$. Portanto, também existem modelos não padrão com
domínio~$\Nat$. Nesses modelos~$\Struct{M}$, naturalmente, pelo menos
alguns números não podem desempenhar seus papéis habituais, pois
algum~$k$ deve ser diferente de~$\Value{\num{n}}{M}$ para todo
$n \in \Nat$.
\end{explain}

\begin{defn}
!!^a{structure}~$\Struct{M}$ para $\Lang{L_A}$ é \emph{computável} se, e somente se,
  $\Domain{M} = \Nat$, $\Assign{\prime}{M}$, $\Assign{+}{M}$ e
  $\Assign{\times}{M}$ são funções computáveis e $\Assign{<}{M}$ é
  uma relação decidível.
\end{defn}

\begin{ex}\ollabel{ex:comp-model-q}
Lembre a estrutura $\Struct{K}$ de \olref[mdq]{ex:model-K-of-Q}.
Seu domínio era $\Domain{K} = \Nat
\cup \{a\}$, com interpretações
\begin{align*}
  \Assign{\Obj{0}}{K} & = 0\\
  \Assign{\prime}{K}(x) & =
  \begin{cases}
    x+1 & \text{se $x\in \Nat$}\\
    a & \text{se $x = a$}
  \end{cases}\\
  \Assign{+}{K}(x, y) & =
  \begin{cases}
    x+y & \text{se $x$, $y \in\Nat$}\\
    a & \text{caso contrário}
  \end{cases}\\
  \Assign{\times}{K}(x, y) & =
  \begin{cases}
    xy & \text{se $x$, $y \in\Nat$}\\
    0 & \text{se $x=0$ ou $y=0$}\\
    a & \text{caso contrário}\\
  \end{cases}\\
  \Assign{<}{K} & =
  \Setabs{\tuple{x,y}}{x, y \in \Nat \text{ e } x<y} \cup
  \Setabs{\tuple{x,a}}{n \in \Domain{K}}
\end{align*}
Mas $\Domain{K}$ é !!{denumerable} e, portanto, tem a mesma cardinalidade
que~$\Nat$. Por exemplo, $g\colon \Nat \to \Domain{K}$, com $g(0) = a$
e $g(n) = n+1$ para $n>0$, é !!a{bijection}. Podemos transformá-la em
um isomorfismo entre um novo modelo~$\Struct{K'}$ de~$\Th{Q}$ e
$\Struct{K}$. Em $\Struct{K'}$, precisamos atribuir funções e relações
diferentes aos símbolos de~$\Lang{L_A}$, pois !!{element}s diferentes
de~$\Nat$ desempenham os papéis de números padrão e não padrão.

Especificamente, $0$ agora desempenha o papel de~$a$, e não o do menor
número padrão. O menor número padrão agora é~$1$. Portanto, atribuímos
$\Assign{\Obj{0}}{K'} = 1$. A função sucessor também é diferente:
dado um número padrão, O.s.\@, um $n > 0$, ela ainda devolve $n+1$.
Mas $0$ agora desempenha o papel de~$a$, que é seu próprio sucessor.
Logo, $\Assign{\prime}{K'}(0) = 0$. Para a adição e a multiplicação,
temos, de modo análogo,
\begin{align*}
\Assign{+}{K'}(x, y) & =
  \begin{cases}
    x+y-1 & \text{se $x$, $y >0$}\\
    0 & \text{caso contrário}
  \end{cases}\\
  \Assign{\times}{K'}(x, y) & =
  \begin{cases}
    1 & \text{se $x = 1$ ou $y = 1$}\\
    xy - x - y + 2 & \text{se $x$, $y > 1$}\\
    0 & \text{caso contrário}\\
  \end{cases}
\end{align*}
E temos $\tuple{x, y} \in \Assign{<}{K'}$ se, e somente se, $x < y$,
$x > 0$ e $y > 0$, ou se $y = 0$.

Todas essas funções são funções computáveis de números naturais, e
$\Assign{<}{K'}$ é uma relação decidível em~$\Nat$---mas elas não são
as mesmas funções sucessor, adição e multiplicação em~$\Nat$, e
$\Assign{<}{K'}$ não é a mesma relação que~$<$ em~$\Nat$.
\end{ex}

\begin{prob}
Dê !!a{structure}~$\Struct{L'}$ com $\Domain{L'} = \Nat$ isomorfa
a~$\Struct{L}$ de \olref[mod][mar][mdq]{ex:model-L-of-Q}.
\end{prob}

\begin{explain}
\Olref{ex:comp-model-q} mostra que $\Th{Q}$ tem modelos não padrão
computáveis com domínio~$\Nat$. No entanto, o resultado a seguir
mostra que isso não vale para modelos de~$\Th{PA}$ (e, portanto,
tampouco para modelos de~$\Th{TA}$).
\end{explain}

\begin{thm}[Teorema de Tennenbaum]
$\Struct{N}$ é o único modelo computável de~$\Th{PA}$.
\end{thm}
  
\end{document}
